\begin{theorem}[连分分母处 \texorpdfstring{$W_1$}{W1} 传输的单箱最优正向置中]\label{thm:xi-kronecker-w1-single-box-optimal-centering}
\leanverified{paper\_xi\_kronecker\_w1\_single\_box\_optimal\_centering}
沿用结论 \ref{con:xi-terminal-kronecker-w1-denominator-one-sided-sensitivity} 的记号。
设 \(p/q\) 为既约有理数且 \(q\ge 2\)，记
\[
\delta:=\alpha-\frac pq.
\]
若满足无跨箱条件
\[
|\delta|<\frac{1}{q(q-1)},
\]
则定理 \ref{thm:kronecker-w1-denominator-closed-form} 的量化耦合计算在该窗口内仍逐字成立，从而
\[
\mathsf{T}_q(\alpha)=
\begin{cases}
\dfrac{1}{2q}+\dfrac{q-1}{2}\left(\dfrac pq-\alpha\right),& \delta<0,\\[2mm]
\dfrac{1}{2q}-\dfrac{q-1}{2}\left(\alpha-\dfrac pq\right)
+\dfrac{q(q-1)(2q-1)}{6}\left(\alpha-\dfrac pq\right)^2,& \delta>0.
\end{cases}
\]
定义
\[
\delta_\ast(q):=\frac{3}{2q(2q-1)}.
\]
则 \(\mathsf{T}_q\) 在整个单箱窗口上的唯一全局极小值恰出现于好侧点
\[
\delta=\delta_\ast(q)>0.
\]
相应最小值为
\[
\boxed{\;
\mathsf{T}_q^{\min}
:=
\mathsf{T}_q\!\left(\frac pq+\delta_\ast(q)\right)
=
\frac1{2q}-\frac{3(q-1)}{8q(2q-1)}.\;}
\]
特别地，
\[
\boxed{\;\mathsf{T}_q^{\min}<\frac1{2q}=\mathsf{T}_q\!\left(\frac pq\right).\;}
\]
因此最优传输置中既不发生在坏侧，也不发生在严格对齐点 \(\alpha=p/q\)，而必须以一个正向箱内偏移实现。
\end{theorem}
\begin{proof}
无跨箱条件给出
\[
(q-1)|\delta|<\frac1q,
\]
故附录定理 \ref{thm:kronecker-w1-denominator-closed-form} 证明中的排序与分位积分计算保持不变，从而同一闭式在该窗口内成立。
当 \(\delta<0\) 时，\(\mathsf{T}_q(\alpha)\) 关于 \(\delta\) 为严格线性函数
\[
\mathsf{T}_q(\alpha)=\frac1{2q}-\frac{q-1}{2}\delta,
\]
故坏侧不存在内点极小值，且其下确界为 \(\delta\uparrow 0\) 时的 \(1/(2q)\)。

当 \(\delta>0\) 时，
\[
\mathsf{T}_q(\alpha)
=
\frac1{2q}-\frac{q-1}{2}\delta
+\frac{q(q-1)(2q-1)}{6}\delta^2,
\]
故
\[
\frac{d^2}{d\delta^2}\mathsf{T}_q(\alpha)=\frac{q(q-1)(2q-1)}3>0,
\]
从而该侧严格凸。唯一临界点由
\[
-\frac{q-1}{2}+\frac{q(q-1)(2q-1)}{3}\delta=0
\]
给出，即
\[
\delta=\delta_\ast(q)=\frac{3}{2q(2q-1)}.
\]
又
\[
\delta_\ast(q)<\frac1{q(q-1)}
\]
等价于 \(3(q-1)<2(2q-1)\)，对一切 \(q\ge 2\) 成立，故该临界点确在窗口内部，并且为好侧唯一极小点。代回得
\[
\mathsf{T}_q^{\min}
=
\frac1{2q}
-\frac{q-1}{2}\cdot\frac{3}{2q(2q-1)}
+\frac{q(q-1)(2q-1)}{6}\cdot\frac{9}{4q^2(2q-1)^2}
=
\frac1{2q}-\frac{3(q-1)}{8q(2q-1)}.
\]
由于减去的第二项严格为正，故 \(\mathsf{T}_q^{\min}<1/(2q)\)。于是全局最优只能发生在好侧正向偏移处。证毕。
\end{proof}

\begin{theorem}[有理阈值处 \texorpdfstring{$W_1$}{W1} 传输曲线的 \texorpdfstring{$C^1$}{C1} 非 \texorpdfstring{$C^2$}{C2} 尖曲率]\label{thm:xi-kronecker-w1-rational-threshold-c1-nonc2-cusp}
\leanverified{paper\_xi\_kronecker\_w1\_rational\_threshold\_c1\_nonc2\_cusp}
沿用定理 \ref{thm:xi-kronecker-w1-single-box-optimal-centering} 的记号。在同一单箱窗口
\[
\left|\alpha-\frac pq\right|<\frac{1}{q(q-1)}
\]
内，函数 \(\alpha\mapsto \mathsf{T}_q(\alpha)\) 在 \(\alpha=p/q\) 处可微，但不属于 \(C^2\)。更准确地，
\[
\partial_-\mathsf{T}_q\!\left(\frac pq\right)
=
\partial_+\mathsf{T}_q\!\left(\frac pq\right)
=
-\frac{q-1}{2},
\]
而左右二阶导数满足
\[
\partial_-^2\mathsf{T}_q\!\left(\frac pq\right)=0,
\qquad
\partial_+^2\mathsf{T}_q\!\left(\frac pq\right)=\frac{q(q-1)(2q-1)}{3}.
\]
因此 \(\alpha=p/q\) 是一条真正的曲率跳跃阈值：一阶几何在两侧无缝拼接，而二阶几何在有理分母处发生不可消去的断裂。
\end{theorem}
\begin{proof}
由定理 \ref{thm:xi-kronecker-w1-single-box-optimal-centering}，当
\(\alpha<p/q\) 时，
\[
\mathsf{T}_q(\alpha)=\frac{1}{2q}+\frac{q-1}{2}\left(\frac pq-\alpha\right),
\]
故左侧导数与二阶导数分别为
\[
\partial_-\mathsf{T}_q\!\left(\frac pq\right)=-\frac{q-1}{2},
\qquad
\partial_-^2\mathsf{T}_q\!\left(\frac pq\right)=0.
\]
当 \(\alpha>p/q\) 时，
\[
\mathsf{T}_q(\alpha)
=
\frac{1}{2q}
-\frac{q-1}{2}\left(\alpha-\frac pq\right)
+\frac{q(q-1)(2q-1)}{6}\left(\alpha-\frac pq\right)^2,
\]
于是
\[
\partial_+\mathsf{T}_q\!\left(\frac pq\right)=-\frac{q-1}{2},
\qquad
\partial_+^2\mathsf{T}_q\!\left(\frac pq\right)=\frac{q(q-1)(2q-1)}{3}.
\]
左右一阶导数相等，故 \(\mathsf{T}_q\) 在 \(\alpha=p/q\) 处属于 \(C^1\)；但左右二阶导数不同，故该点不属于 \(C^2\)。证毕。
\end{proof}

\endinput
